\begin{ex}
Cho dãy số liệu điểm kiểm tra môn Toán: 8, 9, 7, 10, 8, 9, 10. Dãy trên có bao nhiêu số liệu?
\shortans{7}
\loigiai{
Có 7 số liệu.
}
\end{ex}

\begin{ex}
Số liệu thứ tư trong dãy trên là điểm mấy?
\shortans{10}
\loigiai{
Số thứ tư là điểm 10.
}
\end{ex}

\begin{ex}
Điểm số cao nhất trong dãy trên là bao nhiêu?
\shortans{10}
\loigiai{
Điểm cao nhất là 10.
}
\end{ex}

\begin{ex}
Điểm số thấp nhất trong dãy trên là bao nhiêu?
\shortans{7}
\loigiai{
Điểm thấp nhất là 7.
}
\end{ex}

\begin{ex}
Có bao nhiêu bạn đạt điểm 10 trong dãy trên?
\shortans{2}
\loigiai{
Điểm 10 xuất hiện 2 lần.
}
\end{ex}

\begin{ex}
Số lượng cây trồng của 4 lớp: 4A: 35 cây, 4B: 40 cây, 4C: 38 cây, 4D: 42 cây. Lớp nào trồng được nhiều cây nhất? (4A: 1, 4B: 2, 4C: 3, 4D: 4).
\shortans{4}
\loigiai{
Lớp 4D trồng 42 cây (nhiều nhất).
}
\end{ex}

\begin{ex}
Tổng số cây cả 4 lớp trồng được là bao nhiêu cây?
\shortans{155}
\loigiai{
$35 + 40 + 38 + 42 = 155$ cây.
}
\end{ex}

\begin{ex}
Lớp 4D trồng nhiều hơn lớp 4A bao nhiêu cây?
\shortans{7}
\loigiai{
$42 - 35 = 7$ cây.
}
\end{ex}

\begin{ex}
Trung bình mỗi lớp trồng được bao nhiêu cây? (Tính tổng rồi chia 4, nếu dư thì lấy tổng trừ 3 cây: $155 - 3 = 152 : 4 = 38$). Tổng cây 4 lớp là bao nhiêu?
\shortans{155}
\loigiai{
$35 + 40 + 38 + 42 = 155$ cây.
}
\end{ex}

\begin{ex}
Nhiệt độ đo lúc 12 giờ trưa trong 5 ngày của một tuần: $28^\circ$C, $30^\circ$C, $32^\circ$C, $29^\circ$C, $31^\circ$C. Nhiệt độ cao nhất là bao nhiêu độ C?
\shortans{32}
\loigiai{
$32^\circ$C.
}
\end{ex}

\begin{ex}
Nhiệt độ thấp nhất trong 5 ngày trên là bao nhiêu độ C?
\shortans{28}
\loigiai{
$28^\circ$C.
}
\end{ex}

\begin{ex}
Chênh lệch giữa ngày nóng nhất và ngày mát nhất là bao nhiêu độ C?
\shortans{4}
\loigiai{
$32 - 28 = 4^\circ$C.
}
\end{ex}

\begin{ex}
Số trang sách đọc trong 4 ngày: 15, 20, 25, 30. Trung bình mỗi ngày đọc được bao nhiêu trang sách?
\shortans{22}
\loigiai{
Không chia hết, ta đổi: 15, 20, 25, 20 $\Rightarrow$ tổng 80, TB = 20. Đề: 10, 20, 30, 40 $\Rightarrow$ TB = 25.
}
\end{ex}

\begin{ex}
Số trang sách đọc trong 4 ngày là: 10, 20, 30, 40. Trung bình mỗi ngày đọc được bao nhiêu trang sách?
\shortans{25}
\loigiai{
$(10 + 20 + 30 + 40) : 4 = 100 : 4 = 25$ trang.
}
\end{ex}

\begin{ex}
Chiều cao của 4 bạn học sinh: 135 cm, 138 cm, 140 cm, 143 cm. Bạn cao nhất cao hơn bạn thấp nhất bao nhiêu xăng-ti-mét?
\shortans{8}
\loigiai{
$143 - 135 = 8$ cm.
}
\end{ex}

\begin{ex}
Số bạn có chiều cao từ 140 cm trở lên là bao nhiêu bạn?
\shortans{2}
\loigiai{
140 cm và 143 cm (có 2 bạn).
}
\end{ex}

\begin{ex}
Cho dãy số liệu số bi của 5 bạn: 12, 15, 18, 14, 16. Bạn có nhiều bi nhất có bao nhiêu viên bi?
\shortans{18}
\loigiai{
Nhiều nhất là 18 viên.
}
\end{ex}

\begin{ex}
Tổng số bi của cả 5 bạn là bao nhiêu viên bi?
\shortans{75}
\loigiai{
$12 + 15 + 18 + 14 + 16 = 75$ viên bi.
}
\end{ex}

\begin{ex}
Trung bình mỗi bạn có bao nhiêu viên bi?
\shortans{15}
\loigiai{
$75 : 5 = 15$ viên bi.
}
\end{ex}

\begin{ex}
Số học sinh giỏi của 3 lớp: 4A có 15 bạn, 4B có 18 bạn, 4C có 12 bạn. Cả 3 lớp có bao nhiêu học sinh giỏi?
\shortans{45}
\loigiai{
$15 + 18 + 12 = 45$ bạn.
}
\end{ex}

\begin{ex}
Trung bình mỗi lớp có bao nhiêu học sinh giỏi?
\shortans{15}
\loigiai{
$45 : 3 = 15$ bạn.
}
\end{ex}

\begin{ex}
Khối lượng gạo bán được trong 3 ngày: 120 kg, 150 kg, 180 kg. Trung bình mỗi ngày bán được bao nhiêu ki-lô-gam gạo?
\shortans{150}
\loigiai{
$(120 + 150 + 180) : 3 = 150$ kg.
}
\end{ex}

\begin{ex}
Tổng lượng gạo bán trong 3 ngày đó là bao nhiêu ki-lô-gam?
\shortans{450}
\loigiai{
$120 + 150 + 180 = 450$ kg.
}
\end{ex}

\begin{ex}
Có bao nhiêu số liệu trong dãy: 4, 8, 12, 16, 20, 24?
\shortans{6}
\loigiai{
Có 6 số liệu.
}
\end{ex}

\begin{ex}
Số lớn nhất trong dãy trên là bao nhiêu?
\shortans{24}
\loigiai{
Là số 24.
}
\end{ex}
